% ============================================================
%  Chapter 21: The Force-Free Ion --- Partition Depth Minimisation
%  Source: oxford/publications/shader-depth-minimisation/
%          shader-partition-depth-minimisation.tex
% ============================================================

\epigraph{%
  The ion does not follow a force.\\
  It descends the partition depth landscape.%
}{}

\section{The Force-Free Formulation}
\label{sec:ch21-forcefree}

Chapter~\ref{ch:lagrangian} showed that the equation of motion for a bounded
system has the Lorentz-force shape but is force-free: it is gradient descent
on the partition depth field $\Depth(\mathbf{x}, t)$.
Applied to an ion in a mass spectrometer, this means:

\begin{itemize}
  \item There are no ``electric forces''.
    The electric field $\mathbf{E} = -\nabla\Depth$ is the gradient of the
    partition depth field, and the ion moves along this gradient.
  \item There are no ``magnetic forces''.
    The magnetic field appears as the curl of the partition vector potential
    $\Apartition$, and the cyclotron orbit is geodesic motion in the
    $(\Depth, \Apartition)$ landscape.
  \item Fragmentation is not a force event.
    It is a partition-depth transition: the parent ion moves to a new region
    of partition space occupied by its fragments.
\end{itemize}

\section{The Six-Pass GPU Pipeline}
\label{sec:ch21-sixpass}

The partition Lagrangian translates directly into a six-pass GPU rendering
pipeline:

\begin{enumerate}
  \item \textbf{Pass 1 (Ionisation)}: assign initial partition coordinates
    $(n_0, l_0, m_0, s_0)$ from the ESI droplet model.
  \item \textbf{Pass 2 (Drift)}: propagate the ion through the RF multipole
    using the Mathieu stability shader.
  \item \textbf{Pass 3 (Depth descent)}: compute $\nabla\Depth$ in the
    analyser field and advance the ion's partition depth.
  \item \textbf{Pass 4 (Fragmentation)}: detect partition-depth transitions
    above the CID energy threshold; spawn fragment ions.
  \item \textbf{Pass 5 (Detection)}: collect ions at the detector partition
    surface; accumulate image charges.
  \item \textbf{Pass 6 (Spectral reconstruction)}: Fourier-transform the
    image charge time-series to produce $m/z$ and intensity.
\end{enumerate}

Each pass is a fragment shader operating on a texture representing the
partition depth field at that stage.
By the Triple Equivalence (Chapter~\ref{ch:triple}), this computation is
physically equivalent to a measurement: the shader pipeline \emph{is} the
instrument.

\section{Comparison with Four Analyser Technologies}
\label{sec:ch21-comparison}

\begin{table}[htbp]
\centering
\caption{%
  Six-pass GPU pipeline correspondence with physical analyser stages.
}
\begin{tabular}{@{}llll@{}}
\toprule
GPU Pass & Physical Stage & Partition Operation & Observable \\
\midrule
1 & ESI & Depth initialisation & Initial $m/z$ \\
2 & Ion transfer & Mathieu stability filter & Transmission window \\
3 & Mass analysis & Depth gradient descent & $\omega_z$, $\omega_c$, $t_\mathrm{flight}$ \\
4 & CID & Depth transition & Fragment $m/z$ \\
5 & Detection & Boundary accumulation & Image current \\
6 & Spectral output & Koopman FT & $m/z$ spectrum \\
\bottomrule
\end{tabular}
\end{table}

%%  SOURCE: integrate full derivations and figures from
%%          shader-partition-depth-minimisation.tex
%%  ADD: panels 1 (force elimination), 2 (transport/superconductivity),
%%       3 (six-pass pipeline), 4 (GPU journey comparison).

\bigskip
\begin{tcolorbox}[colback=crownblue!5, colframe=crownblue!60!black,
                  title={\textbf{Chapter 21 Summary}}]
\begin{itemize}[noitemsep]
  \item Ion trajectories are force-free: gradient descent on $\Depth(\mathbf{x},t)$.
  \item The six-pass GPU pipeline maps exactly to the six physical stages of a
    mass spectrometric measurement.
  \item By the Triple Equivalence, the GPU pipeline is a physical instrument.
\end{itemize}
\end{tcolorbox}
